\documentclass{ctexbook}
\begin{document}

\chapter{第一章 绪论}
\section{1.1 研究背景与意义}
这是绪论内容，阐述课题的研究背景与研究意义。

\chapter{第二章 过程监测与分析基础方法}
\section{2.1 引言}
第二章属于过程监测与分析基础。在工业制造流程中，各类传感器连续采集动态工艺信号，存在高维与非平稳扰动问题。针对上述技术难题，本章提出动态过程监测与分析框架。本章首先梳理关键工序流程，随后解析变量耦合特性，最后验证表征有效性，为后续章节提供依据。

\section{2.2 过程监测对象特性分析}
本节分析工艺流程中的变量耦合关系与数据特性。

\chapter{第三章 基于自适应门控的状态估计方法}
工业过程数据存在非平稳漂移瓶颈，现场扰动导致瞬态特征难以准确提取，影响产物纯度软测量的响应速度。为此，亟需提高特征提取在多变工况下的自适应表征能力。本章将第 2 章建立的动态校准模型作为基础特征提取器，针对多工况非平稳漂移问题，提出一种自适应注意力门控网络（Adaptive Attention Gating Network, AAGN）。本章首先给出 AAGN 网络的动态门控拓扑与损失约束，随后设计抗噪更新机制，最后在实测数据集上验证表征性能，为第 5 章闭环优化控制提供高保真状态输入。

\section{3.1 自适应门控网络拓扑}
本节给出自适应门控网络的详细结构设计。

\section{3.2 状态估计实验验证}
本节在实测工业数据集上开展对比实验分析。

\chapter{第四章 动态时序鲁棒预测方法}
时序数据缺失导致预测精度下降，在复杂工业运行环境下存在严重的模型鲁棒性不足与泛化难题。第 3 章建立的门控网络在特征缺失时难以保持稳定表征。针对该技术挑战，本章基于变分推断构建鲁棒补偿模型，通过隐空间不确定性建模恢复缺失信息。

\section{4.1 鲁棒预测建模与补偿}
本节给出变分推断鲁棒补偿的具体推导。

\chapter{第五章 多工况闭环协同优化方法}
工业全流程多工况协同优化控制在现代流程工业中具有重要工程应用价值，但实际生产过程中上游工序与下游单元之间存在强烈的动态耦合与大时滞特性，导致全局操作指标难以在单一时间尺度上实现最优寻优，严重影响了综合生产效益与能效水平。为此，亟需建立能够统筹多工况动态响应与多变量交互约束的协同优化控制体系，解决多层级决策之间的信息孤岛瓶颈。第 3、4 章建立的状态估计与鲁棒补偿预测模型为动态工况辨识提供了精确的状态基准，但在遭遇多单元协同负荷阶跃调整时，局部稳态设定值与全流程闭环动态跟踪之间仍存在显著的协调冲突与超调风险。针对上述多工况强耦合协同优化难题，本章基于分层分布式预测控制理论，构建多工况动态闭环协同优化控制方法（Multi-Operating-Condition Closed-Loop Collaborative Optimization, MCCCO），用于实现工序级设定值协同寻优与回路级抗扰跟踪的闭环协同。本方法首先在全流程尺度上构建多目标性能指标函数，将能耗成本与产品合格率转化为多目标 Pareto 寻优问题；其次，设计基于拉格朗日乘子法的双层分布式分解算法，将全局大系统优化问题分解为各工序子系统能够独立并行求解的局部优化子问题，并建立子系统之间的拉格朗日乘子迭代协调通信机制以消除工序间耦合影响；再次，为应对生产负荷切换带来的暂态不平稳扰动，设计带有约束松弛缓冲区的时变预测控制律，在保证物理设备运行安全的前提下实现负荷快速平稳切换；最后，在典型工业生产工况仿真平台与半物理实物测试回路中对所提协同优化控制方法进行多场景对比测试，全面评估该算法在多种典型阶跃扰动与模型失配条件下的闭环响应速度与收敛稳定性，实验结果表明全局能耗得到有效降低，为第 6 章全厂综合工程示范应用提供了关键控制支撑。

\section{5.1 分层协同优化框架设计}
本节详细推导多目标协同优化算法。

\chapter{第六章 工业综合验证与工程应用方法}
第 5 章构建的动态闭环协同优化方法实现了工序级设定值寻优与回路级跟踪的协同，并在半物理测试平台上验证了闭环收敛性，但在面对全厂规模复杂网络通信时延与物理设备硬件限制时，仍存在工程部署与长期运行自适应性不足的局限。实际现场执行机构频繁动作导致阀门机械磨损加剧，影响系统可靠性。

针对上述工程落地难题，本章提出面向工业全流程的长效工程应用方法与自适应保护机制，其核心思想是在协同优化控制回路中引入执行机构平滑死区与通信时延补偿策略。本章首先构建全厂级软硬件系统集成部署架构，随后开展连续 72 小时工业实测运行测试，最后从控制品质、节能降耗以及设备动作频次三方面全面评估工程效益，为全论文研究成果的工业应用推广提供坚实的技术支撑。

\section{6.1 全厂系统集成与部署}
本节介绍工业生产线的软硬件部署环境。

\end{document}
